Recall $f_j(r)$ from Definition~\ref{def:f}. 
The main result of this subsection is the following formula.

\begin{theorem}\label{prop:e_i_formuli}
    For every $s\geq 1$ and $i\geq 0$ we have
    \begin{align*}
        e_{i,(r^s)}
        = \langle \psi^{(r^s)}, \chi^{(1^i) \oplus (rs-i)} \rangle 
        &= \sum_{\gamma\in P_{r,s}(i)} \prod_{j \ge 0}
        (-1)^{(j+1)k_j(\gamma)} \binom{(-1)^{j+1} f_j(r)}{k_j(\gamma)} \\
        &= \sum_{\substack{k_0, \dots, k_r \ge 0 \\ \sum_j k_j = s \\ \sum_j jk_j = i}} \prod_{j = 0}^{r}
        (-1)^{(j+1)k_j} \binom{(-1)^{j+1} f_j(r)}{k_j}.
    \end{align*}
		
    In particular, for $s=1$ we have $e_{i,(r)} = f_i(r)$.
\end{theorem}

\begin{remark}\label{rem:when_factor_is_zero}
    The special case $s=1$ was stated (but not proved) in  Proposition~\ref{prop:e_i_formuli1}.
    The result in that case is not new, as noted in Remark~\ref{rem:e=s}. 
    This case
    shows that $f_i(r) = e_{i,(r)}$ is an inner product of two characters, and is therefore always a non-negative integer. The factor
    \[
        (-1)^{(j+1)k_j} \binom{(-1)^{j+1} f_j}{k_j}
        = \begin{cases}
            \binom{f_j}{k_j}, &\text{ if } j \text{ is odd;}\\
            \binom{f_j+k_j-1}{k_j}, &\text{ if } j \text{ is even}\\
        \end{cases}
    \]
    is therefore also a non-negative integer,
    and is zero if and only if 
    either $j$ is odd and $k_j > f_j$,
    or $j$ is even and $k_j > 0 = f_j$.
    If $k_j = 0$, this factor is equal to $1$ and may be ignored.
\end{remark}
    
In order to prove Theorem~\ref{prop:e_i_formuli} we need a formula for a certain inner product of characters (Lemma~\ref{lem:char_values_and_inner_product}).

First recall some notations from Definition~\ref{def:higher}:
the centralizer $Z_{(r^s)} \cong \ZZ_r \wr S_s$ of an element of cycle type $(r^s)$,
the linear character $\omega^{(r^s)}$ on $Z_{(r^s)}$, 
and the higher Lie character
$\psi^{(r^s)} := \omega^{(r^s)}\uparrow_{Z_{(r^s)}}^{S_n}$. 

Embed $Z_{(r^s)} \cong \ZZ_r \wr S_s$ into 
$K_{r,s} \cong S_r \wr S_s \le S_n$, 
where $\ZZ_r \le S_r$ is generated by a full cycle.
Denote 
\[
    \phi_{r,s} := \omega^{(r^s)}\uparrow_{Z_{(r^s)}}^{K_{r,s}},
\]
so that $\psi^{(r^s)} = \phi_{r,s}\uparrow_{K_{r,s}}^{S_n}$. 
\begin{observation}\label{obs:phi_product} 
    If $s=s_1+s_2$ then $Z_{(r^{s_1})}\times Z_{(r^{s_2})}\leq Z_{(r^s)}$, $K_{r,s_1}\times K_{r,s_2}\leq K_{r,s}$ and also for the characters
    \[ 
        \omega^{(r^s)}
        = \omega^{(r^{s_1})}\otimes\omega^{(r^{s_2})}\text{ and }  \phi_{r,s}\downarrow^{K_{r,s}}_{K_{r,s_1}\times K_{r,s_2}}=\phi_{r,s_1}\otimes\phi_{r,s_2}.
    \]
\end{observation}

In the lemma below we express the multiplicity of a certain linear character in a restriction of $\phi_{r,s}$.
This expression will be used, in the proof of Theorem~\ref{prop:e_i_formuli}, to compute $e_{i,(r^s)}$.
	
For every $0\leq j\leq r$, let 
$R_{r,j} := S_j \times S_{r-j} \le S_r$, in the natural embedding. Then
\[
    R_{r,j} \wr S_s
    = K_{r,s} \cap \left( S_{js} \times S_{(r-j)s} \right).
\]
Denote by $1_{S_n}$ the trivial character and by $\varepsilon_{S_n}$ the sign character of $S_n$, so that
\[
    \chi^{(1^k)\oplus(n-k)}=(\varepsilon_{S_{k}}\times 1_{S_{n-k}})\uparrow_{S_k \times S_{n-k}}^{S_n}.
\]
Define
\[
    \nu_{r,j,s} := (\varepsilon_{S_{js}} \times 1_{S_{(r-j)s}}) \downarrow_{R_{r,j}\wr S_s}^{S_{js} \times S_{(r-j)s}}.
\]
This is a linear character on $R_{r,j} \wr S_s$.
	
\begin{lemma}\label{lem:char_values_and_inner_product}
    For every $0\le j\le r$,
    $f_j(r)$ is a non-negative integer and 
    \[
        \langle \phi_{r,s}\downarrow_{R_{r,j}\wr S_{s}}^{K_{r,s}},\nu_{r,j,s} \rangle 
        = (-1)^{(j+1)s} \binom{(-1)^{j+1} f_j(r)}{s}
        = \begin{cases}
            \binom{f_j(r)}{s}, &\text{if } j \text{ is odd;}\\
            \binom{f_j(r)+s-1}{s}, &\text{if } j \text{ is even.}\\
        \end{cases}
    \]
\end{lemma}
		
\begin{remark}
    As a byproduct, Lemma~\ref{lem:char_values_and_inner_product}
    provides a new proof of the non-negativity of $f_j(r)$.
    Indeed, if $s=1$ then 
    $\langle \phi_{r,1} \downarrow_{R_{r,j}}^{K_{r,1}},\nu_{r,j,1} \rangle = f_j(r)$ 
    is clearly a non-negative integer.
\end{remark}	
	
The rest of this subsection consists of the proofs of Lemma~\ref{lem:char_values_and_inner_product} 
and Theorem~\ref{prop:e_i_formuli}. 
In these proofs 
$r$, $s$ and $j$ are fixed, unless specified otherwise. 
For convenience, we omit the indices and write 
$Z := Z_{(r^s)}$, 
$\omega := \omega^{(r^s)}$, 
$\psi := \psi^{(r^s)}$, 
$K := K_{r,s}$, 
$\phi := \phi_{r,s}$, 
$R := R_{r,j}$, and
$\nu := \nu_{r,j,s}$, 

\begin{proof}[Proof of  Lemma~\ref{lem:char_values_and_inner_product}.]
    The proof consists of two parts. 
    First we determine the character values of the induced character $\phi = \omega\uparrow_Z^K$ on the wreath product $K=S_r\wr S_{s}$; the resulting formula is Equation~\eqref{eq:char_value}. In the second part we 
    apply this formula to compute the inner product.
		
    Let $\zeta: \ZZ_r \to \mathbb{C}$ be the primitive linear character used to define $\omega$; see Definition~\ref{def:higher}. 
    Recall the explicit formula for an induced character \cite[(5.1)]{Isaacs}: 
    For a subgroup $H\leq G$ and a character $\chi$ of $H$, define $\chi^0 : G \to \mathbb{C}$ by $\chi^0(g) = \chi(g)$ if $g\in H$ and $\chi^0(g) = 0$ otherwise. Then
    \begin{equation}\label{eq:induction}
        \chi\uparrow_H^G(y)
        = \frac{1}{|H|} \sum_{x\in G} \chi^0(x^{-1}yx)
        = \sum_{t\in T} \chi^0(t^{-1}yt),
    \end{equation}
    where $T$ is a full set of right coset representatives of $H$ in $G$.
		
    An element of $K=S_r\wr S_s$ can be represented by an $s$-tuple of elements of $S_r$ and a wreathing permutation from $S_s$, so
    $K=\{(x_1,\ldots,x_s;\sigma) \mid x_1,\ldots,x_s \in S_{r},\,\sigma\in S_s\}$
    with the product
    $(x_1,\ldots,x_s;\sigma)(y_1,\ldots,y_s;\tau) = (x_1y_{\sigma^{-1}(1)},\ldots,x_sy_{\sigma^{-1}(s)};\sigma\tau)$. 
    A full set of right coset representatives of $\ZZ_r$ in $S_r$ is $S_{r-1}$ (in the natural embedding).
    Hence, a full set of right coset representatives of 
    $Z = \ZZ_r \wr S_s$ in $K= S_r \wr S_s$ is 
    $T = \{(x_1,\ldots,x_s;1)\mid (\forall i)\, x_i \in S_{r-1}\}$. For any $z_1,\ldots,z_s\in \ZZ_r$, the shifted set $(z_1,\ldots,z_s;1)T$ is also a full set of right coset representatives. 
    Instead of taking the sum over $T$ in \eqref{eq:induction}, we will take it over the union of all the shifted sets, namely $\{(x_1,\ldots,x_s;1) \vert\, (\forall i)\, x_i \in S_r\}$, and divide by $r^s$. 
    Since
    \[
        (x_1,\ldots,x_s;1)^{-1} (y_1,\ldots,y_s;\sigma) (x_1,\ldots,x_s;1)
        = (x_1^{-1}y_1x_{\sigma^{-1}(1)},\ldots,x_s^{-1}y_sx_{\sigma^{-1}(s)};\sigma),
    \]
    we conclude that, 
    for any $y = (y_1,\ldots,y_s;\sigma) \in K$,
    \begin{align*}
        \phi(y) 
        &= \omega\uparrow_Z^K(y_1,\ldots,y_s;\sigma)
        = \frac{1}{r^s} \sum_{x_1,\ldots,x_s \in S_r} \omega^0(x_1^{-1}y_1x_{\sigma^{-1}(1)},\ldots,x_s^{-1}y_sx_{\sigma^{-1}(s)};\sigma) \\ 
        &= \frac{1}{r^s} \sum_{\substack{x_1,\ldots,x_s \in S_r \\ (\forall i)\, x_i^{-1}y_ix_{\sigma^{-1}(i)} \in \ZZ_r}} \omega(x_1^{-1}y_1x_{\sigma^{-1}(1)},\ldots,x_s^{-1}y_sx_{\sigma^{-1}(s)};\sigma) \\
        &= \frac{1}{r^s} \sum_{\substack{x_1,\ldots,x_s \in S_r \\ (\forall i)\, x_i^{-1}y_ix_{\sigma^{-1}(i)} \in \ZZ_r}} \prod_{i=1}^{s} \zeta(x_i^{-1}y_ix_{\sigma^{-1}(i)}).
    \end{align*}

    The factors in the product over $i$ are complex numbers, thus commute. We can therefore rearrange them in an order fitting the decomposition of $\sigma^{-1} \in S_s$ into disjoint cycles: if
    \[
        \sigma^{-1} = C_1 \cdots C_t
    \]
    is a product of $t$ disjoint cycles, choose an element $a_k$ in each cycle $C_k$. Then
    \[
        C_k = (a_k, \sigma^{-1}(a_k), \sigma^{-2}(a_k), \ldots)
        \qquad (1 \le k \le t).
    \]
    Since $\zeta$ is a linear character, cancellation gives
    \begin{align*}
        \prod_{i \in C_k} \zeta(x_i^{-1}y_ix_{\sigma^{-1}(i)})
        &= \zeta(x_{a_k}^{-1}y_{a_k}x_{\sigma^{-1}({a_k})})
        \zeta(x_{\sigma^{-1}({a_k})}^{-1}y_{\sigma^{-1}({a_k})}x_{\sigma^{-2}(a_k)})
        \cdots \\
        &= \zeta(x_{a_k}^{-1}c_kx_{a_k}),
    \end{align*}
    where
    \[
        c_k := y_{a_k} y_{\sigma^{-1}({a_k})} \cdots \in S_r
        \qquad (1 \le k \le t).
    \]
		
    The condition $x_i^{-1}y_ix_{\sigma^{-1}(i)} \in \ZZ_r\, (\forall i)$ implies that the products $x_{a_k}^{-1}c_kx_{a_k}\in \ZZ_r\, (\forall k)$. 
    Hence if $\phi(y) \ne 0$ then, necessarily, each cycle-product $c_k \in S_r$ must be conjugate to an element of $\ZZ_r$.
    Since $\ZZ_r \le S_r$ is generated by a full cycle, 
    a necessary and sufficient condition for $c_k$ to have a conjugate in $\ZZ_r$ is that it is a product of disjoint cycles of the same length.

    For any divisor $d$ of $r$,
    if $c_k \in S_r$ is a product of disjoint $d$-cycles
    and $x_{a_k} \in S_r$ is such that $x_{a_k}^{-1}c_kx_{a_k} \in \ZZ_r$, 
    then the value of $\zeta(x_{a_k}^{-1}c_kx_{a_k})$ is a primitive $d$-th root of unity. By varying the conjugating element $x_{a_k}$, each element of $\ZZ_r$ of order $d$ is obtained with the same multiplicity $|Z_{(d^{r/d})}| = (r/d)! d^{r/d}$.  
    The other $x_i$'s, for $i \in C_k \setminus \{a_k\}$, are arbitrary, as long as $x_i^{-1}y_ix_{\sigma^{-1}(i)}\in \ZZ_r\,(\forall i)$. There are $r^{\ell_k - 1}$ such choices, where $\ell_k$ is the length of the cycle $C_k$.
    We conclude that, for any $y \in K$ for which $c_k \in S_r$ is a product of disjoint $d_k$-cycles $(1 \le k \le t)$,
    \[
        \phi(y) = \frac{1}{r^s} \prod_{k=1}^{t} 
        (r/d_k)! d_k^{r/d_k} r^{\ell_k - 1} 
        \sum_{z \in \ZZ_r \,:\, o(z) = d_k} \zeta(z).
    \]
    If $g$ is a generator of $\ZZ_r$, then $o(g^m) = d$ if and only if $m = jr/d$ for some integer $j$ coprime to $d$. 
    It follows that
    \[
        \sum_{z \in \ZZ_r \,:\, o(z) = d} \zeta(z) 
        = \sum_{0 \le j < d \,:\, (j,d) = 1} \zeta(g^{jr/d})
        = \mu(d).
    \]
    Since $\sum_{k=1}^{t} (\ell_k - 1) = s-t$,
    we now have an explicit formula for the values of $\phi$:
    \begin{equation}\label{eq:char_value}
	   \phi(y) =
	   \begin{cases}\displaystyle
	       \prod_{k=1}^{t} \frac{\mu(d_k) d_k^{r/d_k} (r/d_k)!}{r}, 
	       &\text{if } c_k \text{ is a product of disjoint } d_k \text{-cycles } (\forall k); \\
	       0, &\text{otherwise.}
	   \end{cases}
    \end{equation}

    To determine the inner product $\langle \phi\downarrow_{R\wr S_{s}}^K,\nu\rangle$ we evaluate the linear character~$\nu$ on~$R\wr S_{s}$.
    Let $y = (v_1z_1,\ldots, v_sz_s; \sigma)\in R\wr S_s$, 
    where $v_i\in S_j, z_i\in S_{r-j}$ $(1 \le i \le s)$ and $\sigma\in S_s$. Then, by the definition of $\nu$,
    \[
        \nu(y)
        = \nu(v_1z_1,\ldots, v_sz_s; \sigma)
        = \sgn(\sigma)^j 
        \prod_{i=1}^s \varepsilon(v_i)
    \]
    where $\sgn(\sigma)$ denotes the sign of $\sigma \in S_s$.
    We obtain
    \begin{align*}
        \langle \phi\downarrow_{R\wr S_{s}}^K,\nu\rangle
        &= \frac{1}{|R\wr S_{s}|}
        \sum_{(v_1 z_1,\ldots,v_s z_s;\sigma) \in R\wr S_{s}} 
        \phi(v_1 z_1,\ldots,v_s z_s;\sigma)
        \nu(v_1 z_1,\ldots,v_s z_s;\sigma) \\
        &= \frac{1}{s!} \sum_{\sigma\in S_s} 
        \frac{\sgn(\sigma)^j}{(j!(r-j)!)^s} 
        \sum_{\substack{v_1,\ldots ,v_s \in S_j \\ z_1,\ldots, z_s \in S_{r-j}}}
        \phi(v_1z_1,\ldots,v_sz_s;\sigma)
        \prod_{i=1}^s \varepsilon(v_i).
    \end{align*}
    By Equation~\eqref{eq:char_value},
    for each nonzero summand and each cycle $C_k$ of $\sigma^{-1}$ $(1 \le k \le t)$,
    the cycle-product $c_k \in R \le S_r$ of $y$ has cycle type $d_k^{r/d_k}$, and its restrictions to $S_j$ and $S_{r-j}$ have cycle types $d_k^{j/{d_k}}$ and $d_k^{(r-j)/{d_k}}$, respectively. 
    In particular, $d_k | (r,j)$.
    It follows that the sign
    \[
        \prod_{i=1}^s \varepsilon(v_i)
        = \prod_{k=1}^{t} (-1)^{(d_k+1)j/d_k}. 
    \]
    For any $d | (r,j)$, 
    let $n_d$ denote the number of elements of $R$ which are products of disjoint $d$-cycles. 
    If $C_k$ has length $\ell_k$, then 
    the cycle-product $c_k$ can be a product of disjoint $d_k$-cycles in exactly $(j!(r-j)!)^{\ell_k-1}n_{d_k}$ ways.
    The choices of $d_k$ for different cycles $C_k$  are independent; the only restriction is $d_k|(r,j)$. 
    Using again the equality $\sum_{k=1}^{t} (\ell_k - 1) = s-t$ and Equation~\eqref{eq:char_value}, 
    we obtain 
    \begin{align*}
        \langle \phi\downarrow_{R\wr S_{s}}^K,\nu \rangle
        &= \frac{1}{s!} \sum_{\sigma\in S_s} 
        \frac{\sgn(\sigma)^j}{(j!(r-j)!)^s} 
        \sum_{\substack{v_1,\ldots ,v_s \in S_j \\ z_1,\ldots, z_s \in S_{r-j}}} 
        \phi(v_1z_1,\ldots,v_sz_s;\sigma) 
        \prod_{i=1}^s \varepsilon(v_i) \\
        &= \frac{1}{s!} \sum_{\sigma\in S_s} 
        \frac{\sgn(\sigma)^j}{(j!(r-j)!)^s} 
        \sum_{d_1,\ldots, d_t | (r,j)}\\
        &\, \qquad\qquad\qquad\quad
        \prod_{k=1}^{t} 
        \frac{(j!(r-j)!)^{\ell_k-1} n_{d_k} \mu(d_k) d_k^{r/d_k} (r/d_k)! (-1)^{(d_k+1)j/d_k}}{r} \\
        &= \frac{1}{s!} \sum_{\sigma\in S_s} 
        \frac{\sgn(\sigma)^j}{(j!(r-j)!)^t} 
        \sum_{d_1,\ldots, d_t | (r,j)}
        \prod_{k=1}^{t} 
        \frac{n_{d_k} \mu(d_k) d_k^{r/d_k} (r/d_k)! (-1)^{(d_k+1)j/d_k}}{r} \\
        &= \frac{1}{s!} \sum_{\sigma \in S_s} 
        \frac{\sgn(\sigma)^j}{(j!(r-j)!)^t} 
        \left( \sum_{d|(r,j)} \frac{n_d \mu(d)d^{r/d}(r/d)!  (-1)^{(d+1)j/d}}{r} \right)^t.
    \end{align*}
    Of course, the number of elements of $R = S_j \times S_{r-j}$ which are products of disjoint $d$-cycles is
    \[
        n_d
        = \frac{j!}{d^{j/d} (j/d)!} \cdot \frac{(r-j)!}{d^{(r-j)/d} ((r-j)/d)!}
        = \frac{j!(r-j)!}{d^{r/d}(j/d)!((r-j)/d)!}.
    \]
    Putting everything together
    and recalling that $t = \cyc(\sigma)$ is the number of cycles of $\sigma$,
    we obtain by Definition~\ref{def:f}
    \begin{align*} 
        \langle \phi\downarrow_{R \wr S_{s}}^K,\nu \rangle
        &= \frac{1}{s!} \sum_{\sigma \in S_s} \sgn(\sigma)^j 
        \left( \sum_{d|(r,j)} \frac{\mu(d)(-1)^{(d+1)j/d}}{r} \binom{r/d}{j/d} \right)^{\cyc(\sigma)} \\
        &= \frac{1}{s!} \sum_{\sigma \in S_s} \sgn(\sigma)^j f_j(r)^{\cyc(\sigma)}.
    \end{align*}
    In particular, if $s=1$ then $f = \langle \phi\downarrow_R^K,\nu \rangle$
    is a non-negative integer.
		
    Finally, by~\cite[Proposition 1.3.4]{EC1}, for any $s \ge 0$ and indeterminate $x$, 
    \[
        \frac{1}{s!} \sum_{\sigma \in S_s} x^{\cyc(\sigma)}
        = \binom{x+s-1}{s}.
    \]
    Substituting $-x$ for $x$ and noting that 
    $\sgn(\sigma) = (-1)^{s-\cyc(\sigma)}$, we get
    \[
        \frac{1}{s!} \sum_{\sigma \in S_s} \sgn(\sigma) x^{\cyc(\sigma)} 
        = \binom{x}{s}.
    \]
    This yields the desired formula,
    depending on the parity of $j$,
    for $\langle \phi\downarrow_{R \wr S_{s}}^K,\nu \rangle$.
\end{proof}

To prove Theorem~\ref{prop:e_i_formuli} we need a final ingredient, a combinatorial parametrization of 
$(S_r\wr S_s,S_i\times S_{rs-i})$ double cosets of $S_{rs}$  by partitions.

Recall Definition~\ref{def:partition}. 
The idea and definition of $P_{r,s}(i)$ actually appear
already in the work of Giannelli~\cite{Giannelli}, 
in a similar context but without explicit reference to double cosets or Mackey's formula; 
see Definitions~2.8--2.10 and Proposition~2.11 there.

